\chapter{SQL Level 4: Interview Patterns, SQL for Machine Learning, Performance and Design}\label{ch:sql4}
\companion{techTFQ: Solving real SQL interview queries step by step}{\yts{techtfq+sql+interview+queries+solved+step+by+step}}

This chapter is a catalogue of the patterns that most SQL interview problems are built from. For each: the problem in words, the idea in one sentence,
the query, the real output, and how to explain it aloud. Then the data-scientist-specific part: building leak-free ML feature tables in SQL. Then the
engineering questions every interviewer asks at the end: indexes, query plans, normalisation, OLTP vs OLAP.

\begin{recap}[: how to answer any SQL problem in an interview]
\begin{enumerate}
\item \textbf{Restate the grain}: ``one output row per \emph{what}?'' (per candidate, per day, per department).
\item \textbf{Ask about edge cases}: ties, NULLs, duplicates, days with no data, what to return if nothing qualifies.
\item \textbf{Build in steps} with CTEs; say each step aloud.
\item \textbf{Check with a tiny example}: walk two or three rows through your query.
\item \textbf{Mention performance} only after correctness (indexes on join/filter keys).
\end{enumerate}
\end{recap}

\section{Nth highest value}
\stuck{SQL Nth highest salary: all methods}{\yts{sql+nth+highest+salary+all+methods+dense_rank}}
Salaries (known): 90, 75, 70, 60, 50, 45, 45, 40, 30, 25. Third highest distinct value = 70.
\paragraph{Method 1: DENSE\_RANK (best answer).} Handles ties correctly and generalises to ``per department''.
\runsql{L4_nth_dense}
\paragraph{Method 2: DISTINCT + LIMIT/OFFSET.} Skip $N-1$ distinct values. Returns no row (not NULL) if fewer than $N$ values exist.
\runsql{L4_nth_limit}
\paragraph{Method 3: correlated count (works on any SQL).} A salary is the 3rd highest if exactly 2 distinct salaries are greater.
\runsql{L4_nth_correlated}
\paragraph{Second highest with MAX only.} The largest value below the maximum:
\runsql{L4_second_max}
\begin{trap}
\code{ORDER BY salary DESC LIMIT 1 OFFSET 2} \emph{without} DISTINCT returns the 3rd \emph{row}, which is wrong when the top salaries tie. ROW\_NUMBER has
the same flaw. Also say what happens when there are fewer than $N$ distinct salaries (method 1 returns no row; wrap in
\code{SELECT (subquery) AS nth} to return NULL instead, which is what LeetCode-style graders expect).
\end{trap}

\section{Top N per group}
\stuck{SQL top N per group using window functions}{\yts{sql+top+n+per+group+window+function}}
``Top 2 salaries in each department.'' Rank inside each partition, then filter outside:
\runsql{L4_topn_group}
DataScience returns three people because Kavya and Manav tie for the second-highest salary. If the business wants exactly two people, use ROW\_NUMBER with
a tie-breaker (\code{ORDER BY salary DESC, hire\_date}) and say so. The same pattern gives ``latest application per candidate'', ``most recent salary
per employee'', ``top 3 jobs by applications in each city''.

\section{Duplicates: find, inspect, delete}
\stuck{SQL find and delete duplicate rows}{\yts{sql+find+and+delete+duplicate+rows+row_number}}
Decide first what ``duplicate'' means: same \emph{business key} (candidate, job, date), not the same id.
\runsql{L4_dupes_find}
ROW\_NUMBER labels the copies; every row with \code{rn > 1} is a duplicate to remove:
\runsql{L4_dupes_rownum}
Delete everything except the smallest id in each duplicate group:
\runsql{L4_dupes_delete}
23 rows became 22. In analytics you often do not delete at all: you deduplicate in a CTE (\code{WHERE rn = 1}) at the start of every query, keeping the
raw table untouched.

\section{Self-join classics: employee vs manager}
\runsql{L4_more_than_manager}
Yash (75) earns more than his manager Rahul (70). Variants: ``employees hired before their manager'' (compare hire dates), ``managers with at least 3
reports'' (GROUP BY \code{manager\_id} HAVING \code{COUNT(*) >= 3}: Neha and Arjun, 3 reports each).

\section{Consecutive days: the gaps-and-islands trick}
\stuck{SQL gaps and islands consecutive days problem}{\yts{sql+gaps+and+islands+consecutive+days+problem}}
``Candidates who logged in on at least 3 consecutive days.'' The trick: number each candidate's distinct login dates 1, 2, 3, \dots\ and subtract that
number of days from the date. Consecutive dates move up by one day while the row number also moves up by one, so the difference stays constant; a gap
breaks the run and changes the difference. Each constant value is one \emph{island} (streak).
\runsql{L4_streak_steps}
Aarav's 1--3 March all map to 28 February (one island of length 3); 5--6 March map to 1 March (a new island of length 2). Group by that value and count:
\runsql{L4_streaks}
Note the first CTE: \code{SELECT DISTINCT} removes candidate 109's double login on 4 March. Without it the row numbers would advance twice on one date
and break the islands.

\section{Daily, rolling and cohort metrics}
\stuck{SQL user retention cohort analysis query}{\yts{sql+cohort+retention+analysis+query+tutorial}}
\subsection{Daily active users and rolling active users}
Daily active users with every day shown is the date-spine query from Chapter 47 (\code{L3\_recursive\_dates}). ``Active in the last 3 days'' for each
day uses a range join:
\runsql{L4_rolling_users}
On 3 March, users active on 1--3 March: 101, 102, 104, 105, 109 = 5.

\subsection{Cohort retention}
Group candidates by signup month (their \emph{cohort}) and ask what fraction were active in the first week of March:
\runsql{L4_cohort}
January signups (101--104): 101, 102 and 104 logged in, retention $3/4 = 0.75$. Real retention tables put cohorts on rows and ``months since signup'' on
columns; the logic is the same, with \code{months\_since = month(login) - month(signup)} as the column key.

\subsection{Month-over-month growth}
\runsql{L4_mom}
February: $(11 - 6)/6 = +83.3\%$; March: $(6 - 11)/11 = -45.5\%$. (March is only partial data: say so before anyone panics. Compare like with like, e.g.\
month-to-date vs the same days of last month.)

\section{Funnels and A/B tests}
\stuck{SQL funnel analysis and A/B test analysis}{\yts{sql+funnel+analysis+conversion+rate+query}}
\subsection{Hiring funnel}
Applications move applied $\to$ shortlisted $\to$ interview $\to$ offer. The table stores only the current status, so map each status to the furthest
stage reached and count how many reached at least each stage. (Simplification, said aloud: we treat \code{rejected} as ``stopped at applied'' because we
do not know at which stage the rejection happened. A real events table with one row per stage change would remove this assumption.)
\runsql{L4_funnel}
In SQLite a comparison is 1 or 0, so \code{SUM(reached >= 2)} counts rows; elsewhere write \code{SUM(CASE WHEN reached >= 2 THEN 1 ELSE 0 END)}. 22
applications (duplicate removed), 12 reached shortlisting (54.5\%), 8 interviews, 4 offers: half of interviews convert.

\subsection{A/B test click-through rate}
\runsql{L4_ab_ctr}
Variant B (new ranking) has CTR 0.625 vs 0.25. Is that real? A two-proportion z-test (Chapter 39): pooled $p = 7/16 = 0.4375$, standard error
$\sqrt{p(1-p)(1/8 + 1/8)} = 0.248$, $z = (0.625 - 0.25)/0.248 = 1.51$, two-sided p-value 0.13: \textbf{not significant} with 8 impressions per arm.
SQL computes the counts; the decision needs statistics. Also check the split by day:
\runsql{L4_ctr_daily}
B wins or ties every day, which is reassuring; when a variant wins in every segment but loses overall (or vice versa) you are looking at Simpson's
paradox, caused by unequal segment sizes between arms.

\section{Other patterns worth knowing}
\subsection{Median without a MEDIAN function}
Number the sorted values; the median is the middle row (odd $n$) or the average of the two middle rows (even $n$). With integer division,
\code{(n+1)/2} and \code{(n+2)/2} are the same row when $n$ is odd and the two middle rows when $n$ is even.
\runsql{L4_median}
9 known CTCs; the 5th smallest is 15 (Ananya). PostgreSQL: \code{PERCENTILE\_CONT(0.5) WITHIN GROUP (ORDER BY ctc\_lpa)}.

\subsection{Pivot: rows to columns}
\runsql{L4_pivot}
Conditional aggregation is the portable pivot (pandas: \code{pd.crosstab(jobs.company, jobs.title)}).

\subsection{First event per entity}
\runsql{L4_first_app}

\subsection{Time between two events}
\runsql{L4_time_to_apply}
Candidates apply to ML Engineer postings within a day on average, but take almost 9 days for Data Scientist postings, which could mean DS postings are
less visible or attract more deliberate applicants: a question for the product team, not a conclusion.

\subsection{``Applied to all'' (relational division)}
``Candidates who applied to every Data Analyst job'': count each candidate's distinct DA jobs and compare with the total number of DA jobs.
\runsql{L4_division}
Kabir (103) applied to all three (203, 205, 207).

\subsection{Non-equi join as a data-quality check}
Candidates whose current CTC already exceeds the job's maximum (likely to reject an offer, or a recommendation the ranker should not have made):
\runsql{L4_overqualified}

\section{SQL for machine learning: feature tables without leakage}
\stuck{Feature engineering with SQL for machine learning}{\yts{feature+engineering+with+sql+for+machine+learning}}
In industry, the training table for a model is usually built in SQL: one row per entity (candidate), one column per feature, one label column. Two
mistakes ruin such tables: fan-out and leakage.

\subsection{Fan-out with two one-to-many tables}
Joining candidates to applications \emph{and} logins in one go multiplies rows: candidate 101 has 4 application rows and 5 login rows, so the join creates
$4 \times 5 = 20$ rows and every count is inflated:
\runsql{L4_feature_fanout}
Correct: aggregate each source to one row per candidate in its own CTE, then LEFT JOIN the summaries:
\runsql{L4_feature_table}
Every candidate appears exactly once, including those with no applications or logins (filled with 0), and \code{ctc\_missing} keeps the information
that the fresher had no salary.

\subsection{Leakage: point-in-time correctness}
The feature table above uses \emph{all} history, including events after the moment a prediction would have been made. To train ``will this candidate get
an offer?'', features must use only data before a cutoff date, and the label must come from after it:
\runsql{Q_point_in_time}
Only candidates who had signed up before the cutoff are included (an entity must exist at prediction time). Aarav had 2 applications before 15 February
and got an offer after it (label 1). In production, repeat this for many cutoffs (one snapshot per week) to build a larger training set, and validate on
later cutoffs than you train on (time-based split, Chapter 37).

\begin{trap}[: features that leak]
\code{got\_offer} computed over all time, ``number of interviews'' (which only happens after a shortlist), ``days since last login'' computed at query
time rather than at the cutoff, and global target encodings (offer rate per company over all data) all leak future information. The model looks perfect
offline and fails in production: the same lesson as salary in the placement data (Chapter 42).
\end{trap}

\section{Performance: indexes and query plans}
\stuck{How database indexes work (B-tree) explained}{\yts{how+database+indexes+work+b+tree+explained}}
\subsection{What an index is}
Without an index, finding candidate 105's applications means reading every row (a \emph{full table scan}), like finding a name in an unsorted pile of
forms. An \textbf{index} is a separate, sorted structure (usually a B-tree: a balanced tree whose leaves hold the key values in order with pointers to the
rows), like the alphabetical index at the back of a book. A lookup takes $O(\log n)$ steps instead of $O(n)$: about 30 steps for a billion rows.
\runsql{L4_explain_noindex}
\runsql{L4_explain_index}
\code{SCAN} = read everything; \code{SEARCH ... USING INDEX} = jump straight to the matching rows.

\subsection{Sargable filters}
An index on \code{applied\_date} is sorted by the raw date. Wrapping the column in a function hides that order from the optimiser, so it scans:
\runsql{L4_explain_nonsarg}
Rewrite as a range on the raw column (``\textbf{S}earch \textbf{ARG}ument-\textbf{able}''):
\runsql{L4_explain_sarg}
Other non-sargable forms: \code{WHERE salary * 12 > 600}, \code{LIKE '\%data'} (leading wildcard), \code{WHERE LOWER(email) = ...} (unless there is an index
on \code{LOWER(email)}), implicit type casts (comparing a text column to a number).

\subsection{Rules of thumb}
\begin{itemize}
\item Index columns used in WHERE, JOIN ON, ORDER BY and GROUP BY of frequent queries; foreign keys almost always.
\item A \textbf{composite index} on (\code{cand\_id}, \code{applied\_date}) serves filters on \code{cand\_id} alone or on both, but not on
\code{applied\_date} alone (\emph{leftmost-prefix rule}), like a phone book sorted by surname then first name.
\item A \textbf{covering index} contains every column a query needs, so the table itself is never read.
\item Indexes cost: every INSERT/UPDATE must update them, and they use disk. Low-cardinality columns (gender) rarely benefit.
\item Read the plan (\code{EXPLAIN} / \code{EXPLAIN ANALYZE}) instead of guessing; select only needed columns; filter early; avoid \code{SELECT DISTINCT}
as a fix for fan-out; in warehouses, filter on partition columns (e.g.\ date) so whole partitions are skipped.
\end{itemize}

\section{Schema design: normalisation, OLTP vs OLAP, star schema}
\stuck{Database normalization 1NF 2NF 3NF explained simply}{\yts{database+normalization+1nf+2nf+3nf+explained+simply}}
\subsection{Normalisation in three steps}
Start from one wide sheet: \emph{(app\_id, candidate\_name, candidate\_city, skills, job\_title, company\_name, company\_city, status)}.
\begin{itemize}
\item \textbf{1NF} (atomic values): \code{skills = 'python, sql'} holds a list in one cell; move skills to their own table, one row per (candidate, skill).
\item \textbf{2NF} (no partial dependency): in a table keyed by (cand\_id, job\_id), \code{candidate\_city} depends on cand\_id alone; move it to
\code{candidates}.
\item \textbf{3NF} (no transitive dependency): \code{company\_city} depends on company, which depends on job; move it to \code{companies}.
\end{itemize}
The result is exactly MiniNaukri's design: each fact stored once, so updates cannot create contradictions (\emph{update anomalies}).

\subsection{OLTP vs OLAP}
\begin{center}\small
\begin{tabularx}{\linewidth}{l X X}
\toprule
& OLTP (the product database) & OLAP (the analytics warehouse)\\
\midrule
Workload & Many small reads and writes (apply to a job) & Few large scans and aggregations (applications per city per month)\\
Design & Normalised (3NF), row-oriented & Denormalised star/snowflake schema, column-oriented\\
Examples & MySQL, PostgreSQL & Redshift, BigQuery, Snowflake, Hive, ClickHouse\\
Priorities & ACID, low latency & Throughput, compression, scanning only needed columns\\
\bottomrule
\end{tabularx}
\end{center}
A \textbf{star schema} has a central \emph{fact} table (one row per event: an application, with numeric measures and foreign keys) surrounded by
\emph{dimension} tables (candidate, job, company, date) that describe it. Analysts join the fact to a few dimensions and aggregate. Data scientists usually
query the warehouse, never the production OLTP database.

\section{SQL $\leftrightarrow$ pandas translation}
\stuck{SQL vs pandas: equivalent operations}{\yts{sql+vs+pandas+equivalent+operations+groupby+merge}}
\begin{center}\scriptsize
\begin{tabularx}{\linewidth}{X X}
\toprule
SQL & pandas\\
\midrule
\code{SELECT a, b FROM t WHERE x > 5} & \code{t.loc[t.x > 5, ['a','b']]}\\
\code{ORDER BY a DESC LIMIT 3} & \code{t.sort\_values('a', ascending=False).head(3)} or \code{t.nlargest(3,'a')}\\
\code{SELECT DISTINCT a} & \code{t['a'].drop\_duplicates()} / \code{t['a'].unique()}\\
\code{GROUP BY g} with \code{COUNT(*), AVG(x)} & \code{t.groupby('g').agg(n=('x','size'), avg=('x','mean'))}\\
\code{HAVING COUNT(*) >= 3} & \code{.query('n >= 3')} after the groupby\\
\code{LEFT JOIN u ON t.k = u.k} & \code{t.merge(u, on='k', how='left')}\\
\code{UNION ALL} / \code{UNION} & \code{pd.concat([t, u])} / \code{.drop\_duplicates()}\\
\code{CASE WHEN} & \code{np.where}, \code{np.select}, \code{pd.cut}\\
\code{COALESCE(x, 0)} & \code{t.x.fillna(0)}\\
\code{AVG(x) OVER (PARTITION BY g)} & \code{t.groupby('g').x.transform('mean')}\\
\code{ROW\_NUMBER() OVER (PARTITION BY g ORDER BY d)} & \code{t.sort\_values('d').groupby('g').cumcount() + 1}\\
\code{DENSE\_RANK() OVER (ORDER BY x DESC)} & \code{t.x.rank(method='dense', ascending=False)}\\
\code{LAG(x) OVER (PARTITION BY g ORDER BY d)} & \code{t.sort\_values('d').groupby('g').x.shift()}\\
\code{SUM(x) OVER (ORDER BY d)} & \code{t.sort\_values('d').x.cumsum()}\\
Pivot with \code{SUM(CASE ...)} & \code{pd.pivot\_table} / \code{pd.crosstab}\\
\bottomrule
\end{tabularx}
\end{center}
One semantic difference to mention: pandas' \code{groupby} drops NaN keys by default (\code{dropna=False} keeps them), whereas SQL's GROUP BY keeps NULL as
its own group.

% ======================================================================
\section{Interview questions: Level 4}
\stuck{Advanced SQL interview questions for data scientists}{\yts{advanced+sql+interview+questions+data+scientist}}

\begin{qa}{\classic Find the second (and the Nth) highest salary. What if there are ties or fewer than N salaries?}
\oneline{\code{DENSE\_RANK() OVER (ORDER BY salary DESC)} and keep rank N; ties are handled because DENSE\_RANK ranks distinct values; if fewer than N
distinct values exist, return NULL by wrapping the query as a scalar subquery.}
\why Ranking distinct values is exactly what ``Nth highest salary'' means. \code{MAX(salary) WHERE salary < (SELECT MAX(salary) ...)} works for N = 2 only.
\code{LIMIT 1 OFFSET N-1} needs DISTINCT. The correlated version (count distinct greater values = N-1) works on any SQL, at $O(n^2)$ cost.
\worked Here: second highest = 75 (\code{L4\_second\_max}); third = 70 by all three methods (\code{L4\_nth\_dense}, \code{L4\_nth\_limit},
\code{L4\_nth\_correlated}).
\begin{lstlisting}[style=sql]
SELECT (SELECT DISTINCT salary FROM
          (SELECT salary, DENSE_RANK() OVER (ORDER BY salary DESC) AS dr
           FROM employees WHERE salary IS NOT NULL) t
        WHERE dr = 20) AS nth_salary;   -- returns one row with NULL
\end{lstlisting}
\pushes \emph{Per department?} Add \code{PARTITION BY dept}. \emph{Without window functions or LIMIT?} The correlated count. \emph{What about the NULL
salary?} Exclude it explicitly; in a DESC sort SQLite puts NULL last, PostgreSQL puts NULL first, and RANK would then give NULL rank 1 in PostgreSQL.
\pitfall ROW\_NUMBER or LIMIT without DISTINCT.
\end{qa}

\begin{qa}{\classic Find and delete duplicate rows, keeping one copy.}
\oneline{Define the business key, find duplicates with \code{GROUP BY key HAVING COUNT(*) > 1}, label copies with \code{ROW\_NUMBER() OVER (PARTITION BY
key ORDER BY id)}, and delete rows with \code{rn > 1} (or keep \code{MIN(id)} per key).}
\why Duplicates almost never share the surrogate id; they share the business key (candidate, job, date). ROW\_NUMBER's ORDER BY decides which copy is kept
(the first-inserted, the latest-updated, the most complete).
\worked \code{L4\_dupes\_find} finds (101, 201, 2026-01-16) twice; \code{L4\_dupes\_delete} keeps \code{MIN(app\_id)} and leaves 22 rows.
In PostgreSQL:
\begin{lstlisting}[style=sql]
DELETE FROM applications WHERE app_id IN (
  SELECT app_id FROM (
    SELECT app_id, ROW_NUMBER() OVER (PARTITION BY cand_id, job_id, applied_date
                                      ORDER BY app_id) AS rn
    FROM applications) t
  WHERE rn > 1);
\end{lstlisting}
\pushes \emph{How do you prevent it?} A UNIQUE constraint on the business key, or idempotent inserts (\code{INSERT ... ON CONFLICT DO NOTHING}).
\emph{Table has no id at all?} Use the database's row id (\code{ctid}, \code{rowid}) or rebuild with \code{CREATE TABLE ... AS SELECT DISTINCT}.
\pitfall Deleting all copies (\code{WHERE key IN (duplicates)}), including the one to keep.
\end{qa}

\begin{qa}{\classic Find users who logged in on at least 3 consecutive days.}
\oneline{Deduplicate (user, date), subtract ROW\_NUMBER (in days) from the date to get a constant per run, group by user and that constant, keep groups
with count $\ge 3$.}
\why Within a run of consecutive dates, both the date and the row number increase by 1 per row, so their difference is constant; any gap increases the
date by more than the row number and starts a new group.
\worked \code{L4\_streak\_steps} shows Aarav's dates mapping to 28 Feb, 28 Feb, 28 Feb, 1 Mar, 1 Mar; \code{L4\_streaks} returns 104 (4 days), 101 (3)
and 107 (3).
\pushes \emph{Alternative?} LAG: flag rows where \code{date - LAG(date) > 1} as a new streak, then a running SUM of the flags gives the streak id.
\emph{Longest streak per user?} \code{MAX(streak\_len)} per user over the grouped result. \emph{Consecutive weeks?} Use week numbers instead of dates.
\pitfall Forgetting to remove same-day duplicates.
\end{qa}

\begin{qa}{\classic Compute 7-day retention: of the users who signed up on each day, what fraction were active 7 days later?}
\oneline{Cohort table (user, signup\_date) LEFT JOIN activity on \code{user = user AND activity\_date = signup\_date + 7}, group by signup date,
\code{COUNT(activity.user) / COUNT(*)}.}
\why The LEFT JOIN keeps users who were not retained (they count in the denominator); the join condition defines ``retained''. Variants: ``active on day
7'' (exact) or ``active any time in days 1--7'' (\code{BETWEEN}).
\worked The monthly-cohort version on MiniNaukri, \code{L4\_cohort}: January 0.75, February 0.5, March 0.5.
\begin{lstlisting}[style=sql]
SELECT c.signup_date, COUNT(*) AS cohort,
       COUNT(DISTINCT l.cand_id) AS retained_d7,
       ROUND(1.0 * COUNT(DISTINCT l.cand_id) / COUNT(*), 2) AS d7_retention
FROM candidates c
LEFT JOIN logins l ON l.cand_id = c.cand_id
                  AND l.login_date = date(c.signup_date, '+7 days')
GROUP BY c.signup_date;
\end{lstlisting}
On MiniNaukri this returns 0 for every signup day, because the login log only covers 1--6 March and no candidate's day 7 falls in it: a reminder to check
that the activity data actually covers each cohort's window.
\pushes \emph{Why might the latest cohorts look bad?} They have not had 7 days yet (right-censoring); exclude cohorts younger than the window.
\emph{Why COUNT(DISTINCT) in the numerator?} A user can log in twice on day 7.
\pitfall Putting the activity condition in WHERE (drops non-retained users, giving 100\% retention).
\end{qa}

\begin{qa}{\classic Top 3 jobs by number of applications in each city.}
\oneline{Count applications per job (LEFT JOIN so zero-application jobs count), rank with \code{DENSE\_RANK() OVER (PARTITION BY city ORDER BY n DESC)},
keep rank $\le 3$.}
\worked
\begin{lstlisting}[style=sql]
WITH n AS (
  SELECT j.job_id, j.city, COUNT(a.app_id) AS n_apps
  FROM jobs j LEFT JOIN applications a ON a.job_id = j.job_id
  GROUP BY j.job_id, j.city
), r AS (
  SELECT *, DENSE_RANK() OVER (PARTITION BY city ORDER BY n_apps DESC) AS dr FROM n
)
SELECT city, job_id, n_apps FROM r WHERE dr <= 3 ORDER BY city, n_apps DESC;
\end{lstlisting}
\why Grain: one row per job first, then rank within city. Bengaluru has 4 jobs (201: 5 including the duplicate, 207: 3, 206: 2, 202: 1), so 202 is cut.
\pushes \emph{Ties?} State whether you want DENSE\_RANK (all ties) or ROW\_NUMBER (exactly 3).
\pitfall Ranking before aggregating (ranking individual application rows).
\end{qa}

\begin{qa}{Build a hiring funnel from an \code{application\_events} table (one row per status change). How is that better than a status column?}
\oneline{Count distinct applications that reached each stage (\code{COUNT(DISTINCT CASE WHEN stage = 'interview' THEN app\_id END)}), divide consecutive
stages; an events table records the path, so ``rejected after interview'' is distinguishable from ``rejected at screening''.}
\why A single status column overwrites history; MiniNaukri's funnel had to assume rejected applications stopped at stage 1 (\code{L4\_funnel}). With events
you also get time between stages (median days to shortlist) and where in the funnel rejections happen.
\worked With the current-status approximation: 22 applied, 12 shortlisted (54.5\%), 8 interviews, 4 offers (50\% of interviews).
\pushes \emph{Ordered funnel (must pass stage 2 after stage 1)?} Compare per-application timestamps of each stage with MIN() per stage and require
increasing times.
\pitfall Counting events instead of distinct applications (a stage can be logged twice).
\end{qa}

\begin{qa}{Your query computing applications per day is slow on a 2-billion-row table. How do you investigate?}
\oneline{Read the plan with EXPLAIN (ANALYZE); check that the date filter is sargable and hits a partition or index; select only needed columns;
pre-aggregate (materialised daily table); check statistics and join strategy.}
\why Typical culprits: \code{WHERE DATE(applied\_at) = ...} (function on the column defeats indexes and partition pruning), \code{SELECT *} on a columnar
warehouse (reads every column), joining before filtering, fan-out joins, \code{COUNT(DISTINCT)} on huge cardinalities (approximate with
\code{APPROX\_COUNT\_DISTINCT}/HyperLogLog when exactness is not needed).
\worked \code{L4\_explain\_nonsarg} shows \code{SCAN} for \code{strftime('\%Y', applied\_date) = '2026'}; the half-open range in \code{L4\_explain\_sarg}
uses the index.
\pushes \emph{Composite index for ``applications of candidate X in March''?} (\code{cand\_id}, \code{applied\_date}): equality column first, range
column second. \emph{Why not index everything?} Writes slow down and storage grows.
\pitfall Guessing optimisations without looking at the plan.
\end{qa}

\begin{qa}{What is normalisation? When would you deliberately denormalise?}
\oneline{Organising tables so each fact is stored once (1NF atomic values, 2NF no partial dependency on part of a key, 3NF no transitive dependency),
which prevents update anomalies; denormalise in analytics warehouses and read-heavy paths, where fewer joins matter more than update cost.}
\why A normalised OLTP schema makes writes safe and small. Analytical queries then need many joins; a warehouse star schema (wide fact + dimensions) or a
flat feature table trades redundancy for speed and simplicity. Updates there come from batch pipelines, not users, so anomalies are controlled.
\worked The single wide application sheet with \code{skills = 'python, sql'} and company city repeated per row becomes MiniNaukri's 4 tables plus a
skills table.
\pushes \emph{What is BCNF?} A stricter 3NF: every determinant is a candidate key. \emph{Star vs snowflake?} Snowflake normalises the dimensions further
(company $\to$ industry table).
\pitfall Claiming normalised is always better.
\end{qa}

\begin{qa}{\deep You are asked for ``average applications per active candidate per day''. Write down three different numbers this could mean, and the SQL
for the one you would pick.}
\oneline{(1) total applications / total candidate-days with activity; (2) the mean over days of (applications that day / active candidates that day);
(3) the mean over candidates of their own daily average; they differ because of weighting, and the choice must be agreed with the stakeholder.}
\why (1) weights busy days more; (2) gives each day equal weight (an average of ratios); (3) gives each candidate equal weight, so a candidate active one
day counts as much as one active every day. ``Active'' also needs a definition: logged in, or applied? Each choice is a different metric.
\worked Definition (2), with active = logged in:
\begin{lstlisting}[style=sql]
WITH a AS (SELECT applied_date AS d, COUNT(*) AS apps FROM applications GROUP BY 1),
     u AS (SELECT login_date AS d, COUNT(DISTINCT cand_id) AS users FROM logins GROUP BY 1)
SELECT AVG(1.0 * COALESCE(a.apps, 0) / u.users) AS avg_apps_per_active_user
FROM u LEFT JOIN a ON a.d = u.d;
\end{lstlisting}
\pushes \emph{Which would you report?} The one matching the decision: capacity planning wants (1); a product health dashboard usually wants (2) with
the numerator and denominator shown. \emph{Why show both parts?} So that a change in the ratio can be traced to a change in applications or in users.
\pitfall Writing SQL before defining the metric.
\end{qa}

\begin{qa}{\deep You build a training table in SQL for ``will this candidate get an offer in the next 30 days?'' and the offline AUC is 0.99. What do you
check first?}
\oneline{Leakage: every feature must be computed only from data before the prediction cutoff, the label only from after it, and the population must be
the candidates who existed (and were eligible) at the cutoff.}
\why Features computed over all history (``number of interviews'', ``has offer'', ``days since last login'' computed today) contain the answer. Global
aggregates (company offer rate over all time) leak too. A random train/test split of snapshots also leaks, because the same candidate's later snapshots
inform earlier ones; split by time.
\worked \code{Q\_point\_in\_time}: features from applications strictly before 15 February, label from on or after it, only candidates who signed up
before the cutoff. Aarav: 2 earlier applications, label 1 (offer on 19 February).
\pushes \emph{How do feature stores help?} They store timestamped feature values and do ``as-of'' joins automatically. \emph{Sanity checks?} Feature
importance dominated by one suspicious feature; performance collapsing on a later time-split; features whose values change for the past when recomputed.
\pitfall Celebrating the AUC.
\end{qa}

\begin{qa}{\deep Your A/B test query shows variant B with CTR 0.625 vs A with 0.25. The product manager wants to ship B today. What do you say?}
\oneline{The counts are tiny (8 impressions per arm): a two-proportion z-test gives $z = 1.51$, p = 0.13, so the difference is not significant; we also
need to check randomisation balance, the unit of analysis (users, not impressions) and novelty effects before shipping.}
\why SQL gives counts; inference needs the sampling variability. Impressions from the same user are correlated, so the effective sample size is the number of
users and the variance should be computed per user (or by the delta method). A sample-ratio-mismatch check (are arms the expected size?) catches broken
assignment. Novelty: users click new things at first.
\worked Pooled $p = 7/16$; SE $= \sqrt{0.4375 \times 0.5625 \times (1/8 + 1/8)} = 0.248$; $z = 0.375/0.248 = 1.51$; two-sided $p = 0.13$.
\pushes \emph{How many impressions would you need?} For 80\% power to detect 0.25 $\to$ 0.35 at $\alpha = 0.05$: about $n \approx (1.96 + 0.84)^2
[0.25\cdot0.75 + 0.35\cdot0.65]/0.1^2 \approx 325$ per arm. \emph{Daily results?} \code{L4\_ctr\_daily}: B is never worse on any day, encouraging but
not proof.
\pitfall Treating a SQL ratio as a conclusion.
\end{qa}

\begin{qa}{\deep Why is \code{WHERE YEAR(applied\_at) = 2026} slower than \code{WHERE applied\_at >= '2026-01-01' AND applied\_at < '2027-01-01'}
when there is an index on \code{applied\_at}?}
\oneline{The index stores raw \code{applied\_at} values in sorted order; the optimiser can binary-search a range on the raw column, but it cannot search
for values of a function of the column without computing the function for every row, so it scans.}
\why A B-tree supports lookups on its key's order. \code{YEAR(x) = 2026} is equivalent to a range, but the optimiser does not invert arbitrary functions.
Expression indexes (\code{CREATE INDEX ... ON t (YEAR(applied\_at))}) or partitioning on the year are alternatives.
\worked SQLite plans: \code{SCAN applications} for the strftime filter vs \code{SEARCH applications USING INDEX idx\_app\_date (applied\_date>? AND
applied\_date<?)} for the range (\code{L4\_explain\_nonsarg}, \code{L4\_explain\_sarg}).
\pushes \emph{Same issue elsewhere?} \code{WHERE amount + 0 = 5}, \code{WHERE CAST(id AS TEXT) = '5'}, \code{LIKE '\%x'}, and comparing a VARCHAR column
with an integer literal (implicit cast on the column).
\pitfall Answering ``functions are slow''; the real reason is losing the index order.
\end{qa}
